\section{Computational assertions and their audit scope}
\label{sec:computational}

The classification in this working draft is a logical consequence of the
analytic arguments and the six assertions below.  We state the assertions
as mathematical claims, separately from the evidence currently available
for them.  Their names do not upgrade the status of an archived calculation.
Several source packages retain the designation
\texttt{DIAGNOSTIC\_ONLY / NOT\_EVIDENCE}; a project decision to use a result
does not itself constitute an independent numerical proof.  The repositories
are public, and the fixed source identifiers are provided in the references.
An independently implemented replay of every required enclosure is not
claimed in this revision.

Throughout this section, $E$ abbreviates the reference-coordinate energy
$\widehat E_\lambda$ and every spatial derivative is taken in reference
coordinates.  Put
\[
 \mathcal D=\{(\rho,z):\rho,z\geq0,\ \rho^2+z^2\leq1\},\qquad
 H=\frac{E_\rho}{\rho},\qquad V=\frac{E_z}{z}.
\]
At interior points with $\rho=0$ or $z=0$, the respective quotients
are extended by the corresponding second derivative.  Every inequality
for a second or higher spatial derivative is an assertion on the
intersection of the stated box with the open ball, even when that box
meets the boundary.  At the boundary, first derivatives mean the
continuous inner trace of the full gradient; $H$ is then used only
where $\rho>0$, as that radial trace divided by $\rho$, and $V$ only
where $z>0$.  No boundary Hessian extension is assumed in this section.
The full-gradient trace condition is distinct from a constrained
critical-point condition on the boundary.

\subsection{The finite statements}

\begin{computationalassertion}[Finite oblate transverse and axial signs]
\label{comp:c1}
For $1/32\leq\lambda\leq1$ and interior $(\rho,z)\in\mathcal D$,
\[
 H(\rho,z;\lambda)>0.
\]
The same strict sign holds for the radial trace quotient at boundary
points with $\rho>0$.
For $0<z\leq1$, the axial first derivative satisfies
\[
 \begin{array}{ll}
 E_z(0,z;\lambda)<0,&1/32\leq\lambda\leq2/5,\\
 E_z(0,z;\lambda)>0,&33/50\leq\lambda\leq1.
 \end{array}
\]
On the axis segment $0\leq z\leq1/8$, one has, more specifically,
\[
 \begin{array}{ll}
 E_{zz}(0,z;\lambda)<0,&1/32\leq\lambda\leq2/5,\\
 E_{zz}(0,z;\lambda)>0,&33/50\leq\lambda\leq1.
 \end{array}
\]
The corresponding first-derivative signs follow by integrating
from $z=0$; on $1/8\leq z\leq1$, $E_z$ is enclosed directly.
\end{computationalassertion}

The archived finite cover uses either $H$ itself or $X=E_\rho$.
An $X$ leaf is accepted only when its entire box has $\rho>0$ and
its lower endpoint for $X$ is positive.  Thus its implication for $H$
has a stated denominator condition.  A box is discarded as empty only
after its intersection with $\mathcal D$ has been shown to be empty.
The overlap of the parameter slabs is intentional: $2/5$ is covered
through an upward dyadic endpoint, and the upper family begins at
$21/32<33/50$.
On $[21/32,1]$, recorded lower endpoints for $H$, axial $E_z$
on $z\geq1/8$, and axial $E_{zz}$ on $z\leq1/8$ are respectively
\[
 \frac{11239578}{2^{36}},\qquad
 \frac{28921220}{2^{36}},\qquad
 \frac{34030169745}{2^{36}}.
\]
Assertion~\ref{comp:c1} excludes every off-axis critical point in this
finite oblate band; the axis census in the remaining window is the
separate scalar statement below.  It makes no assertion about a uniform
transverse sign for $0<\lambda<1/32$.
The latter range is handled by the analytic thin-oblate argument \citep{data-oblate-full}.

\begin{computationalassertion}[Oblate scalar axis census]
\label{comp:c2}
Let
\[
 g(t,\lambda)=E_z(0,t;\lambda),\quad
 h_o(\lambda)=g_t(0,\lambda),\quad
 c_3(\lambda)=\frac16 g_{ttt}(0,\lambda),\quad
 b(\lambda)=g(1,\lambda).
\]
The following statements hold.
\begin{enumerate}
 \item On $[1/4,1]$, $h_o'(\lambda)>0$, and
 $h_o(2/5)<0<h_o(83/200)$.  Its unique zero $c_o$ therefore satisfies
 \[
  \frac25<c_o<\frac{83}{200}.
 \]
 Moreover, $c_3(\lambda)<0$ throughout $[2/5,83/200]$.
 \item The polar trace has a unique zero $b_o$ in $[2/5,33/50]$,
 is negative before that zero and positive after it, and
 \[
  \frac{674802}{2^{20}}<b_o<\frac{674811}{2^{20}}.
 \]
 The recorded endpoint inequalities are
 \[
  b(674802/2^{20})\leq-\frac{175746}{2^{36}},\qquad
  b(674811/2^{20})\geq\frac{153930}{2^{36}}.
 \]
 \item On $[2/5,33/50]$, the positive interior zeros of $g$ obey
 \[
  \#\{t\in(0,1):g(t,\lambda)=0\}
  =\begin{cases}
    1,&c_o<\lambda<b_o,\\
    0,&2/5\leq\lambda\leq c_o\ \text{or}\ b_o\leq\lambda\leq33/50.
   \end{cases}
 \]
 At every such nonzero interior root, $g_t<0$.
 At $\lambda=b_o$, the endpoint $t=1$ is the only positive zero
 in $(0,1]$; it is not counted as an interior zero.
\end{enumerate}
\end{computationalassertion}

The scalar root count is assembled from specified sign gates, not from
samples of a plotted branch.  In the lower half, write
$g(t,\lambda)=t\Phi(t^2,\lambda)$.
On $[2/5,9/20]\times[0,1/2]$, the archived fourth-energy-derivative
gate $g_{ttt}<0$, together with the analytic transfer identity,
gives $\partial_\tau\Phi<0$.  Signed values at $t=1/2$ and the center
then determine the lower-half roots.  The upper-half cover on
$[2/5,9/20]$ has 160 parameter slabs, split into 48 root-free slabs,
78 slabs with endpoint signs and a negative $g_t$ strip, and
34 slabs with an isolated root tube.  On $[9/20,5/8]$, 159 accepted
parameter intervals have positive left-side values, $g_t<0$ in the
root tube, and negative right-side values or a negative polar trace.
There are 124 tubes ending before the pole and 35 ending at the pole.
The residual outer region for the former 124 is separately covered.
On $[5/8,33/50]\times[31/32,1]$, the gates are $g_t<0$,
$g(31/32,\lambda)>0$, and the signed polar trace.  This description
specifies which inequalities provide existence, uniqueness, and exclusion;
the complete closed-interval assembly and its source receipts are archived
in the axis reference \citep{data-oblate-axis}.

To make the logical use of C2(3) explicit, let $\lambda_x$ be the
unique zero of $g(1/2,\lambda)$ in
$(699/1600,281/640)$ (the C1a gate).  The following is the
\emph{gate-to-root-count assembly} from the fixed axial record
\citep{data-oblate-axis}.  Here $\Phi(t^2,\lambda)=g(t,\lambda)/t$
for $t>0$, $B_o(\lambda)=g(1,\lambda)$, ``tube'' denotes an interval
with opposite endpoint signs and $g_t<0$, and $N_+$ counts zeros in
$(0,1)$ only.  The entries are conditional on the specified sign gates;
they do not independently validate their interval implementations.

\begin{center}
\begingroup\footnotesize
\setlength{\tabcolsep}{3pt}\renewcommand{\arraystretch}{1.12}
\begin{tabularx}{\linewidth}{@{}p{29mm}p{43mm}Xr@{}}
\toprule
$\lambda$ & Lower axis $0<t\le1/2$ & Upper axis $1/2<t<1$ and outer exclusion & $N_+$\\
\midrule
$[2/5,\cO)$ & $\Phi(0)<0$, $\Phi_\tau<0$: no root & Gap-1 L: $g<0$ & $0$\\
$\{\cO\}$ & $\Phi(0)=0$, $\Phi_\tau<0$: no positive root & Gap-1 L: $g<0$ & $0$\\
$(\cO,83/200]$ & $\Phi(0)>0>\Phi(1/4)$, strictly decreasing: one root & Gap-1 L: $g<0$ & $1$\\
$(83/200,\lambda_x)$ & C1a/A1, $\Phi(0)>0>\Phi(1/4)$: one root & Gap-1 X: upper interval excluded & $1$\\
$\{\lambda_x\}$ & $\Phi(1/4)=0$: only $t=1/2$ & X: all $t>1/2$ excluded & $1$\\
$(\lambda_x,9/20)$ & $\Phi(1/4)>0$: no lower root & X/R: one monotone upper root, negative outer region & $1$\\
$[9/20,5/8]$ & C1c: $g>0$ & C1b: 159 accepted tubes; 124 outer regions and Gap-2 excluded & $1$\\
$(5/8,\bO)$ & C1d: $g>0$ through $31/32$ & Band: $g_t<0$, $g(31/32)>0>B_o$ & $1$\\
$\{\bO\}$ & C1d: $g>0$ through $31/32$ & Band: $g_t<0$; only $t=1$ is a zero & $0$\\
$(\bO,33/50]$ & C1d: $g>0$ through $31/32$ & Band: $g_t<0$, $B_o>0$: no upper root & $0$\\
\bottomrule
\end{tabularx}
\endgroup
\end{center}
The endpoint $83/200$ is assigned to the third row; Gap-1's slab
ownership of that value is X rather than L and supplies a compatible
sign check. The X/R junction at $703/1600$ is a shared certified
root, not two roots. The value $9/20$ belongs to C1b, and $5/8$ to
the C1b row; the C1d/band check at $5/8$ counts that same root
once. The cut $t=31/32$ is bookkeeping, while $t=1$ is a boundary
zero and is never included in $N_+$. The fixed source also records
the 48 L, 78 X and 34 R parameter slabs and the 124 unclamped and
35 clamped C1b tubes. The present ZIP includes a separately readable
C2 evidence map; full C2 source-code replay remains a separate task.

The prior centre computation used 160-bit Arb followed by a 192-bit replay of the transcribed kernel \citep{data-center-axis}. For a separate endpoint test write $s=\sin t$, $m=\cos t$, $D_0=s^2+\lambda^2 m^2$ and $\theta=\arctan((\lambda-1/\lambda)ms)$. On the axis the signed angle is $\arctan(\lambda(m-z)/s)-\arctan(m/(\lambda s))$. Differentiate its square with weight $1-zm$ twice at $z=0$ and use reflection to obtain
\[
 h_o(\lambda)=\int_0^{\pi/2}
 \frac{s\,[2\lambda^2s^2+4\lambda(1-\lambda^2)m s^3\theta]}{D_0^2}\,dt .
\]
At the endpoints of the finer interval in Table~\ref{tab:critical}, complex ball integration gives an upper bound $-186781876279625303/2^{96}$ on the left and a lower bound $10585972787056429/2^{96}$ on the right. Combined with the monotonicity in C2 this isolates the same zero. The wider interval in C2 also supports its cubic-coefficient test.

\begin{computationalassertion}[The near-spherical band]
\label{comp:c3}
For $1\leq\lambda\leq2$ and $(\rho,z)\in\mathcal D$,
\begin{align*}
 E_z(\rho,z;\lambda)&\geq\frac{16189}{2^{36}}
       &&(1/4\leq z\leq1),\\
 E_{zz}(\rho,z;\lambda)&\geq\frac{66330426}{2^{36}}
       &&(0\leq z\leq1/4).
\end{align*}
On $z=0$ one has $E_{yyyy}<0$ throughout $0\leq\rho<1$. On the boundary equator,
\[
 H_b(\lambda):=E_\rho(1,0;\lambda)
 \geq\frac{32725738}{2^{36}}>0.
\]
\end{computationalassertion}

Reflection gives $E_z(\rho,0)=0$, so the first two inequalities imply
$V\geq16189/2^{36}$ whenever $z\ne0$.
The whole-equator sign and positive boundary value exclude all noncentral equatorial points. The earlier 15-leaf equatorial test instead enclosed either $H\geq1383839082/2^{36}$ or $E_{\rho\rho}\leq-4138135153/2^{36}$ at positive radius. It gives a separate one-way-crossing argument.
The source's correction E2 withdraws the earlier claim that the first
two terms $4/3-8r^2/9$ of the spherical expansion are an exact formula.
None of the inequalities in Assertion~\ref{comp:c3} uses that truncation \citep{data-near-sphere}.

\begin{computationalassertion}[The finite prolate window]
\label{comp:c4}
For $2\leq\lambda\leq5$ and $(\rho,z)\in\mathcal D$,
\begin{align*}
 E_{zz}&\geq\frac{1503705}{2^{33}}
          &&(0\leq z\leq1/4),\\
 E_z&\geq\frac{3341881}{2^{36}}
          &&(1/4\leq z\leq1).
\end{align*}
For $0\leq\rho<1$ on the equator, $E_{yyyy}\leq-109964349/2^{36}$. At each positive-radius interior zero of $H$, the separate shape-derivative input gives $H_\lambda<0$.
For $0\leq\rho\leq1/8$ on the equator,
$E_{yyyy}\leq-5355762265/2^{35}$, where $y=q_2$.
The center coefficient $H_c(\lambda)=H(0,0;\lambda)$ satisfies
\[
 H_c'(\lambda)\leq-\frac{1090349669}{2^{34}}
 \qquad(2\leq\lambda\leq5).
\]
The boundary coefficient satisfies
\[
 \begin{array}{ll}
 H_b'(\lambda)\leq-6439538363/2^{36},&2\leq\lambda\leq9/4,\\
 H_b(\lambda)\leq-11303125/2^{36},&9/4\leq\lambda\leq5.
 \end{array}
\]
The unique boundary and center zeros, respectively $b_p$ and $c_p$,
are enclosed by
\[
 b_p\in\left(\frac{34651360}{2^{24}},\frac{34651369}{2^{24}}\right),
 \qquad
 c_p\in\left(\frac{79261996}{2^{24}},\frac{79262005}{2^{24}}\right).
\]
At the displayed left and right endpoints the bounds are, respectively,
\[
 \begin{array}{c|cc}
 &\text{left endpoint, lower bound}&\text{right endpoint, upper bound}\\ \hline
 H_b&6485992/2^{47}&-5189437/2^{47}\\
 H_c&3088214/2^{47}&-2555623/2^{47}.
 \end{array}
\]
On $[2,9/4]\times[7/8,1)\times\{z=0\}$,
$E_{\rho\rho}\leq-28849345377/2^{36}$ uniformly up to the
boundary equator.  Only the first-derivative trace is used at
$\rho=1$.
\end{computationalassertion}

The first two inequalities give $V\geq3341881/2^{36}$ off the
equator.  The fourth-derivative bound gives $H_\rho<0$ for every $0<\rho<1$ by \eqref{eq:fourth_derivative}.
A separate parameter-direction cover has alternatives $H>0$,
$H<0$, and $H_\lambda<0$.
Thus every root satisfies both strict derivative inequalities,
and implicit differentiation gives
$\rho_*'(\lambda)=-H_\lambda/H_\rho<0$.
The quartic gate controls the degenerate center, while the endpoint
signs locate the two transitions \citep{data-window}.

\begin{computationalassertion}[The finite part of the prolate tail]
\label{comp:c5}
Let $l=\sqrt{\lambda^2-1}$ and use the one-dimensional integrals
$N_c,N_e,N_p$ of~\eqref{eq:pt-three-coefficient-densities}.
Set
\[
 A_c(l)=3\pi/4-N_c(l),\quad
 A_e(l)=3\pi/4-N_e(l),\quad
 A_p(l)=\pi/2-N_p(l).
\]
The anchor inequalities are
\[
 \begin{array}{c|ccc}
 \lambda&A_c&A_e&A_p\\ \hline
 5&<1&<1&<-3/16\\
 8&<-1/20&<-2/5&<-3/5.
 \end{array}
\]
In addition, on the closed domain
\[
 5\leq\lambda\leq8,\qquad(\rho,z)\in\mathcal D,\qquad
 \sqrt{1-z^2}\geq1/\lambda,
\]
the transverse coefficient has the following upper bounds:
\[
 \begin{array}{c|c}
 \lambda\text{ interval}&H\text{ upper bound}\\ \hline
 {[5,6]}&-68778309/2^{36}\\
 {[6,7]}&-9210649/2^{33}\\
 {[7,8]}&-1095717/2^{30}.
 \end{array}
\]
\end{computationalassertion}

These statements concern finite integrals and finite covers only.
The extension to all $\lambda\geq5$, including the complementary
polar region, follows from the analytic comparison of the coefficients
in Section~\ref{sec:prolate-tail}.  In particular, a numerical check
at $\lambda=8$ is not being extrapolated without that argument.
The complete exact anchor endpoints are retained in the accompanying
source notes \citep{data-chord}.

\begin{computationalassertion}[The polar lower-envelope integral]
\label{comp:c6}
For the explicitly defined lower-bound integral
$\mathcal B$ in~\eqref{eq:pt-cap-computational-integral},
\[
 \mathcal B(\epsilon,w,u)\geq
 \frac{47473218194072764542240096942353646083}{2^{127}}
 >\frac14
 \quad\text{on}\quad
 [0,1/5]\times[0,1]\times[0,1].
\]
\end{computationalassertion}

Here $z=\sqrt{1-\epsilon^2w^2}$ and $\rho=\epsilon wu$.
For $\epsilon>0$, the analytic derivation identifies
$\lambda=1/\epsilon$ and proves $\lambda E_z\geq\mathcal B$.
The endpoint $\epsilon=0$ is a continuous extension of this
lower-bound integral, not a definition of the energy at an infinite
parameter.  The analytic subtraction terms in $\mathcal B$ account
for the omitted angular regions.  Assertion~\ref{comp:c6} does not
assume a global asymptotic remainder estimate \citep{data-cap}.

\subsection{Arithmetic, integration, and singularity handling}
\label{subsec:computational-methods}

\paragraph{Directed binary64 enclosures.}
The finite oblate, near-sphere, and prolate-window kernels use a
custom interval type with IEEE binary64 endpoints.  The program sets
\texttt{FE\_UPWARD}; lower endpoints of arithmetic operations are
obtained by the appropriate sign reversal, and upper endpoints by
upward rounding.  A positive square root has an upward endpoint
and a lower endpoint obtained by stepping that result down one
representable value; zero is handled exactly.  The recorded build flags
include \texttt{-frounding-math}, \texttt{-fno-fast-math}, and
\texttt{-ffp-contract=off}.  Their correctness depends on the stated
floating-point environment and compiler behavior, not merely on
printing many digits.  A startup directed-rounding check is present.

The usual coefficient tables have 8,192 cells and degree-60 positive
series with an outward tail allowance $2^{-38}$ for the coefficient
quantities and derivatives used by those kernels.
Dyadic angular grids are constructed through square-root half-angle
identities.  Outer-box coordinates are exact integers in units of
$2^{-20}$; the ordinary reported integral endpoints use units of
$2^{-36}$.  For the oblate polar refinement, the identities
$J'=Q$ and $K'=gQ$ give mean-value enclosures from a fixed table node.

The principal integral enclosure on a spatial cell is the range of
the complete kernel on that cell, multiplied by its positive measure.
Summing these products is a Darboux enclosure; it does not omit an
unestimated quadrature remainder.  The graded oblate chart
$m=1-(1-t)^2$ uses exact cell differences in $m$ as weights.
Its expression $C=(1-m^2)+\lambda^2m^2$ preserves positivity near
the thin poles.  Known geometric ratio bounds are used on contact
cells instead of dividing an interval that contains zero.

\paragraph{Fourth-order remainder for the critical-value refinement.}
The prolate-window location calculations use a separate tensor
Simpson enclosure.  Their coefficient tables have 4,096 cells,
degree 96, and normalized derivative tails through order six bounded
by $2^{-52}$.  On a panel of width $\Delta$, the Hermite remainder
has the sign-controlled average enclosure
\[
 \operatorname{mean}(f)-\operatorname{Simpson}(f)
 \in -\frac{\Delta^4}{120}\,[f^{(4)}/4!].
\]
Indeed, the normalized nodal polynomial is
$t(t-1)(t-1/2)^2\leq0$ on $[0,1]$ and its integral is $-1/120$.
The code encloses the fourth Taylor coefficient on each whole panel.
In two variables, the two integration errors are decomposed in
sequence with positive outer weights.  The critical endpoint
enclosures are then reported outward in units of $2^{-47}$.
This error control is separate from an ordinary point quadrature
or a comparison between two numerical resolutions.

\paragraph{Arb coefficients and an integer finite-tail kernel.}
The six one-dimensional anchors needed in Assertion~\ref{comp:c5}
are enclosed with 96-bit Arb, 2,048 uniform $m$ cells, and the
degree-12 lower polynomial for $Q$ in
\eqref{eq:pt-positive-series}.  Its omitted coefficients are
nonnegative, so discarding them gives the required lower bound.
The independent audit uses 4,096 cells and degree 14.
The finite $[5,8]$ bridge uses fixed-point interval endpoints at
scale $2^{-24}$ with signed 128-bit integer intermediates.
Square roots are corrected until the defining integer-square
inequalities hold; a floating-point initial guess is not accepted
without that correction.  Its 2,048-cell coefficient table is
generated by 144-bit Arb, with degree 16 and the explicit positive
tail bound for $P$ in~\eqref{eq:pt-series-enclosures}; $Q$ is bounded
below by its positive truncation.  Exact trigonometric moments are
enclosed in the table.  Keeping the two negative chord contributions
paired preserves the identity
\eqref{eq:pt-paired-negative-identity}.

\paragraph{The fixed polar cover.}
Assertion~\ref{comp:c6} uses 128-bit Arb and
$16\times8\times8=1,024$ exact outer boxes.  The retained
$\theta$ interval is divided into 96 equal cells, and the
$\phi$ interval into 48.  Every entire-cell range is multiplied
by its area.  On this retained angular rectangle, the positive
lower bounds for the denominators are proved before the computation;
the two omitted regions are bounded analytically.

\paragraph{Boundary contact.}
The direct gradient kernel is bounded at all noncontact surface points
and admits a dominated inner limit.  A contact point has area zero.
For the Hessian enclosures, the chord change of variables cancels
the vanishing endpoint distance before evaluation.  In particular,
the bounded chord expression is extended at a collapsing chord;
the general singular expression is not numerically evaluated there.
The boundary-equator derivative also has a dedicated kernel.
Higher spatial derivatives used for the degenerate center are evaluated
away from the physical boundary.

\subsection{Which finite computations have been checked?}
\label{subsec:computational-workload}
\begin{center}\small
\begin{tabularx}{\linewidth}{lXX}\toprule
Input & Computation supporting the classification & Scope of the additional checks\\\midrule
C1 & Oblate exclusion and axial signs & 123 earlier selected items. New Arb check of the weakest positive axial-gradient leaf among 277 in the upper oblate band.\\
C2 & Central and polar zeros, and axial root count & New endpoint computations for $c_o$ at separation $10^{-12}$. The prior root-count and monotonicity records supply uniqueness.\\
C3 & Vertical exclusion and positive boundary trace & 35 earlier selected leaves. New Arb check of the weakest among 28,358 positive vertical-gradient leaves. A second implementation re-encloses all 176 additional equatorial boxes.\\
C4 & Vertical exclusion, radial monotonicity and motion of the circle & All 731 radial boxes checked with a second implementation. Earlier separate checks include 22 vertical leaves and a selected shape-derivative leaf.\\
C5 & Finite bridge and six tail anchors & All six anchors recomputed in Arb. Refinement of 79 bridge parents uses the existing chord kernel.\\
C6 & Polar integral & All 1,024 outer boxes recomputed with a different integration implementation and angular partition.\\\bottomrule
\end{tabularx}\end{center}
The supplement identifies the source populations and the distinction between a replay, a second implementation of an integral, and a separately derived formula. The full vertical and shape-direction populations have different audit coverage from the radial population.
